\documentclass[10pt,a4paper]{amsart}
\usepackage[margin=19mm]{geometry}
\usepackage{amsmath,amssymb,mathtools}
\usepackage[scr=rsfs,cal=euler]{mathalfa}
\usepackage{xfrac}
\usepackage[unicode,hidelinks,bookmarks=false]{hyperref}
\newcommand{\ie}{i.e.\ }
\newcommand{\cf}{cf.\ }
\newcommand{\resp}{respectively\ }
\newcommand{\Subsection}{\subsection}
\providecommand{\nobreakdash}{\nobreak\hbox{-}}
\setlength{\emergencystretch}{2em}
\multlinegap0pt
\usepackage{amssymb}
\usepackage{dsfont}
\newtheorem{lemma}{Lemma}
\newtheorem{proposition}{Proposition}
\newtheorem{theorem}{Theorem}
\newcommand{\I}{\mathcal I}
\newcommand{\Span}{\operatorname{span}}
\newcommand{\one}{\mathbf 1}
\title{Functional identities of degree 2 at two-sided zero products on incidence algebras}
\author{Hongyu Jia, Zhankui Xiao}
\date{Source: arXiv:2608.26853v1 (CC BY 4.0)}
\begin{document}
\maketitle

\noindent\emph{Source and licence.} Functional identities of degree 2 at two-sided zero products on incidence algebras. Version arXiv:2608.26853v1 (CC BY 4.0), distributed by arXiv under the Creative Commons Attribution 4.0 International licence, http://creativecommons.org/licenses/by/4.0/, as recorded in the arXiv licence metadata for this exact version and shown on the abstract page https://arxiv.org/abs/2608.26853v1. Full licence notice: \texttt{licenses/arxiv-2608.26853-CC-BY-4.0.txt}.\par\smallskip

\noindent\textit{Editorial scope.} Counted unit: the article's theorem thm:main-characterization (source0.tex lines 1361--1384). The article's complete original proof of it is delivered with it (source0.tex lines 1386--1579, one \texttt{proof} environment of 194 lines). No mathematics is written, repaired or re-derived: the statement and the proof are the article's own lines, in the article's own notation, and no retained line is substituted or re-flowed.  Retained uncounted as the source-local context the proof needs: lines 246--255; lines 365--373; lines 459--462; lines 718--720; lines 990--999; lines 1015--1017; lines 1329--1339.  Nothing was omitted from any retained span, so the package carries no omission record.  No retained line contains a citation, so the wrapper carries no bibliography and no bibliography entry.  No retained line contains a figure environment, an image-inclusion command or drawing code, so no image is delivered and the package declares no asset dependency.  Editorial formatting only, in addition: an AMS \texttt{amsart} preamble that restates the article's own declarations and macros used by the retained lines, namely the environments \texttt{lemma}, \texttt{proposition}, \texttt{theorem} on separate counters, together with the standard packages those lines need.  The retained environments print in this document as Lemma 1, Proposition 1, Theorem 1 in order of appearance, whereas the article prints the same items with the numbering of its own full document.  The complete native source of the article as distributed in the arXiv e-print for this exact version is bundled unchanged as \texttt{sources/208668/source0.tex}.
\begin{lemma}\label{lem:graph-equivalences}
Let $X$ be a finite connected poset with $|X|>2$. The following statements are
equivalent:
\begin{enumerate}
\item[(a)] Any two edges in $E(\Gamma(X))$ are contained in one cycle.
\item[(b)] $X$ is $2$-connected.
\item[(c)] $\Gamma(X)\setminus\{z\}$ is connected for every $z\in X$.
\item[(d)] $\Gamma(X)$ has no cut vertex.
\end{enumerate}
\end{lemma}

\begin{proposition}\label{prop:center}
Let $X_1,\ldots,X_m$ be the connected components of
$\Gamma(X)$ and set $\one_i=\sum_{x\in X_i}e_x$.
Then the center
\[
Z(I(X,R))=\bigoplus_{i=1}^m R\one_i.
\]
In particular, if $X$ is connected, then $Z(I(X,R))=R\one$.
\end{proposition}

\begin{lemma}\label{lem:derivation-diagonal}
If $D:\I\to\I$ is a derivation, then
$D(e_x)(y,y)=0$ for all $x,y\in X$.
\end{lemma}

\begin{equation}\label{eq:antisymmetric-local}
A(f)g-A(g)f+fB(g)-gB(f)=0,~\text{whenever}~fg=gf=0.
\end{equation}

\begin{proposition}\label{thm:antisymmetric-standard-form}
Assume that $|X|>2$ and any two edges in $E(\Gamma(X))$ are contained in one cycle.
Then the $R$-linear maps $A,B:\I\to\I$ satisfy \eqref{eq:antisymmetric-local} if and only if there exist $C_1,C_2\in\I$, a derivation
$D:\I\to\I$, and a linear central-valued map
$\Omega:\I\to Z(\I)$ such that
\begin{equation}\label{eq:antisymmetric-standard-form}
A(f)=C_1f+D(f)+\Omega(f)~\text{and}~
B(f)=fC_2+D(f)+\Omega(f).
\end{equation}
\end{proposition}

\begin{equation}\label{eq:symmetric-local}
G(f)g+fH(g)+G(g)f+gH(f)=0,~\text{whenever}~fg=gf=0.
\end{equation}

\begin{proposition}\label{thm:symmetric-standard-form}
Assume that $|X|>2$ and any two edges in $E(\Gamma(X))$ are contained in one cycle.
The the $R$-linear maps $G,H:\I\to\I$ satisfy \eqref{eq:symmetric-local} if and only if
there exist $C_1,C_2\in\I$, a derivation
$D:\I\to\I$, and a linear central-valued map
$\Omega:\I\to Z(\I)$ such that
\begin{equation}\label{eq:symmetric-standard-form}
G(f)=C_1f+D(f)+\Omega(f)~\text{and}~
H(f)=fC_2+D(f)-\Omega(f).
\end{equation}
\end{proposition}

\begin{theorem}\label{thm:main-characterization}
Let $R$ be a commutative ring with unity such that $\frac{1}{2}\in R$.
Let $X$ be a connected finite poset with $|X|>2$ and set $\I=I(X,R)$.
The following statements are equivalent:
\begin{enumerate}
\item[(i)] Any two edges in $E(\Gamma(X))$ are contained in one cycle.
\item[(ii)] Let $F_1,F_2,F_3,F_4:\I\to\I$ be $R$-linear maps satisfying
\begin{equation}\label{eq:four-map-local}
F_1(f)g+fF_2(g)+F_3(g)f+gF_4(f)=0,
\end{equation}
whenever $fg=gf=0$. Then there exist
$C_1,C_2,C_3,C_4\in\I$, derivations $D_1,D_2:\I\to\I$, and linear
central-valued maps $\Omega_1,\Omega_2:\I\to Z(\I)$ such that
\begin{equation}\label{eq:four-map-standard}
\begin{aligned}
F_1(f)&=C_1f+D_1(f)+\Omega_1(f),\\
F_2(f)&=fC_2+D_1(f)+\Omega_2(f),\\
F_3(f)&=C_3f+D_2(f)-\Omega_2(f),\\
F_4(f)&=fC_4+D_2(f)-\Omega_1(f),
\end{aligned}
\end{equation}
for all $f\in\I$.
\end{enumerate}
\end{theorem}

\begin{proof}
Assume that (i) holds. Let $F_1,F_2,F_3,F_4$ be $R$-linear maps satisfying
\eqref{eq:four-map-local}. Since the condition $fg=gf=0$ is
symmetric in $f$ and $g$, interchanging them gives
\begin{equation}\label{eq:four-map-swapped}
F_1(g)f+gF_2(f)+F_3(f)g+fF_4(g)=0.
\end{equation}
Subtracting \eqref{eq:four-map-swapped} from \eqref{eq:four-map-local} gives
$$0=(F_1-F_3)(f)g-(F_1-F_3)(g)f+f(F_2-F_4)(g)-g(F_2-F_4)(f).$$
Applying Proposition~\ref{thm:antisymmetric-standard-form} to
$A=F_1-F_3$ and $B=F_2-F_4$, there are $P_1,P_2\in\I$, a derivation $\Delta_1$, and a
central-valued map $\Lambda_1$ such that, for all $f\in\I$,
\begin{equation}\label{eq:difference-forms}
\begin{aligned}
(F_1-F_3)(f)&=P_1f+\Delta_1(f)+\Lambda_1(f),\\
(F_2-F_4)(f)&=fP_2+\Delta_1(f)+\Lambda_1(f).
\end{aligned}
\end{equation}
Adding \eqref{eq:four-map-swapped} to \eqref{eq:four-map-local} gives
$$0=(F_1+F_3)(f)g+f(F_2+F_4)(g)+(F_1+F_3)(g)f+g(F_2+F_4)(f).$$
Applying Proposition~\ref{thm:symmetric-standard-form} to
$G=F_1+F_3$ and $H=F_2+F_4$, there are $P_3,P_4\in\I$, a derivation $\Delta_2$, and a
central-valued map $\Lambda_2$ such that
\begin{equation}\label{eq:sum-forms}
\begin{aligned}
(F_1+F_3)(f)&=P_3f+\Delta_2(f)+\Lambda_2(f),\\
(F_2+F_4)(f)&=fP_4+\Delta_2(f)-\Lambda_2(f).
\end{aligned}
\end{equation}
Set
\[
\begin{aligned}
C_1&=\tfrac12(P_1+P_3),&
C_2&=\tfrac12(P_2+P_4),&
C_3&=\tfrac12(P_3-P_1),&
C_4&=\tfrac12(P_4-P_2),\\
D_1&=\tfrac12(\Delta_1+\Delta_2),&
D_2&=\tfrac12(\Delta_2-\Delta_1),\\
\Omega_1&=\tfrac12(\Lambda_1+\Lambda_2),&
\Omega_2&=\tfrac12(\Lambda_1-\Lambda_2).
\end{aligned}
\]
Then it is clear that $D_1,D_2$ are derivations and $\Omega_1,\Omega_2$ are central-valued maps.
Combining \eqref{eq:difference-forms} and \eqref{eq:sum-forms}, we have
\[
\begin{aligned}
F_1(f)
&=\frac12\bigl((F_1-F_3)(f)+(F_1+F_3)(f)\bigr)=C_1f+D_1(f)+\Omega_1(f),\\
F_2(f)
&=\frac12\bigl((F_2-F_4)(f)+(F_2+F_4)(f)\bigr)=fC_2+D_1(f)+\Omega_2(f),\\
F_3(f)
&=\frac12\bigl((F_1+F_3)(f)-(F_1-F_3)(f)\bigr)=C_3f+D_2(f)-\Omega_2(f),\\
F_4(f)
&=\frac12\bigl((F_2+F_4)(f)-(F_2-F_4)(f)\bigr)=fC_4+D_2(f)-\Omega_1(f),
\end{aligned}
\]
for all $f\in\I$. This proves \eqref{eq:four-map-standard} and we get (ii).

Now assume that (ii) holds. We assume opposite that there are two edges in $E(\Gamma(X))$ that are not contained in a cycle.
Then $\Gamma(X)$ has a cut vertex $z$ by Lemma~\ref{lem:graph-equivalences}.
Let $X_1,\ldots,X_m$ be the connected components of $X\setminus\{z\}$, where $m\geq2$.
Set $p_i=\sum_{x\in X_i}e_x$ and let
\begin{equation*}
\mathcal{J}_i
=R\text{-}\Span\{e_{uv}\mid u\leq v,\ u,v\in X_i\cup\{z\},\,
(u,v)\neq(z,z)\}.
\end{equation*}
We next establish some facts about these $R$-submodules $\mathcal{J}_i$ of $\I$.

\smallskip
\noindent\emph{Claim 1.}
If $i\neq j$, then
\begin{equation}\label{eq:Ji-cross-zero}
\mathcal J_i\mathcal J_j=0=\mathcal J_j\mathcal J_i.
\end{equation}

We only prove $\mathcal J_i\mathcal J_j=0$. Suppose, to the contrary, that $\mathcal J_i\mathcal J_j\neq0$.
Then there exist $e_{xy}\in\mathcal J_i$ and $e_{uv}\in\mathcal J_j$ such that
$e_{xy}e_{uv}\neq 0$, which implies $y=u$. Notice that $y\in X_i\cup\{z\}$ and $u\in X_j\cup\{z\}$, but
$X_i$ and $X_j$ are distinct connected components of $X\setminus\{z\}$.
Therefore, $(X_i\cup\{z\})\cap(X_j\cup\{z\})=\{z\}$ and hence $y=u=z$. On the other hand,
by the definition $(x,y)\neq (z,z)$, but $y=z$, we can get that $x\in X_i$.
Similarly, there is $v\in X_j$. That means $x,z,v$ are three distinct vertices and $x<z<v$.
Then $x<v$ and hence the vertices $x$ and $v$ are
adjacent in the comparability graph $\Gamma(X)$. Moreover, the edge $\{x,v\}$ belongs to
$\Gamma(X)\setminus\{z\}$. This edge connects a vertex of $X_i$ to
a vertex of $X_j$, contradicting the fact that $X_i$ and $X_j$ are
distinct connected components of $\Gamma(X)\setminus\{z\}$.
Therefore $\mathcal J_i\mathcal J_j=0$.

\smallskip
\noindent\emph{Claim 2.}
For all $h\in\mathcal J_i$ with $1\leq i\leq m$,
\begin{equation}\label{eq:local-unit}
(p_i+e_z)h=h=h(p_i+e_z).
\end{equation}

In fact, since $X_i\cap \{z\}=\emptyset$, we have
\[
(p_i+e_z)e_{uv}=\left(\sum_{x\in X_i}e_x+e_z\right)e_{uv}=e_{uv},
\]
for any basis element $e_{uv}\in \mathcal J_i$. Therefore, we obtain $(p_i+e_z)h=h$, and
similarly $h=h(p_i+e_z)$, for all $h\in\mathcal J_i$ with $1\leq i\leq m$.

\smallskip
We next construct an $R$-linear map which is Lie derivable at zero.
Notice that every basis element $e_{xy}$ of $\I$, except $e_z$, belongs to exactly one $\mathcal{J}_i$.
Hence there is a direct sum decomposition of $R$-modules as follows
\begin{equation}\label{eq:I-cut-decomposition}
\I=Re_z\oplus\bigoplus_{i=1}^m\mathcal J_i.
\end{equation}
Let $t_1,\ldots,t_m\in R$ be scalars, which are not all equal. We define an $R$-linear map $L:\I\to\I$ by
\[
L(e_z)=-\sum_{i=1}^m t_ip_i
\]
and
\[
L(e_{uv})=t_i e_{uv}
\]
for all $e_{uv}\in \mathcal{J}_i$ and $1\leq i\leq m$. For arbitrary elements $f,g\in \I$,
we write them according to \eqref{eq:I-cut-decomposition} as
$f=\alpha e_z+\sum_{i=1}^m f_i$ and $g=\beta e_z+\sum_{i=1}^m g_i$, where
$f_i,g_i\in\mathcal J_i$ and $\alpha,\beta\in R$.
Then it follows from the definition of $L$ that
\begin{equation}\label{eq:L-cut-formula}
L(f)=\sum_{i=1}^m t_i(f_i-\alpha p_i)~\text{and}~
L(g)=\sum_{i=1}^m t_i(g_i-\beta p_i).
\end{equation}
By Claim~1, a direct computation shows
\begin{equation}\label{eq:commutator-cut-components}
[f,g]=\sum_{i=1}^m h_i,
\end{equation}
where
$h_i=[f_i,g_i]+\alpha[e_z,g_i]+\beta[f_i,e_z]\in\mathcal J_i$.
Using \eqref{eq:L-cut-formula} and Claim~1, we also have
\[
\begin{aligned}
{}[L(f),g]
&=\sum_{i=1}^m t_i
[f_i-\alpha p_i,\,\beta e_z+g_i]\\
&=\sum_{i=1}^m t_i
\bigl([f_i,g_i]+\beta[f_i,e_z]-\alpha[p_i,g_i]\bigr)
\end{aligned}
\]
and
\[
\begin{aligned}
{}[f,L(g)]
&=\sum_{i=1}^m t_i
[\alpha e_z+f_i,\,g_i-\beta p_i]\\
&=\sum_{i=1}^m t_i
\bigl([f_i,g_i]+\alpha[e_z,g_i]-\beta[f_i,p_i]\bigr).
\end{aligned}
\]
Notice that the Claim~2 implies $[p_i+e_z,h]=0$, for any $h\in\mathcal J_i$. Hence
$[p_i,g_i]=-[e_z,g_i]$ and $[f_i,p_i]=-[f_i,e_z]$. Therefore,
\begin{equation}\label{eq:L-cut-commuting-calculation}
[L(f),g]+[f,L(g)]=2\sum_{i=1}^m t_i h_i.
\end{equation}
If we assume $[f,g]=0$, then it follows from \eqref{eq:commutator-cut-components} and \eqref{eq:I-cut-decomposition}
that $h_i=0$ for each $1\leq i\leq m$. Hence $L$ is Lie derivable at zero by \eqref{eq:L-cut-commuting-calculation}, i.e.,
\begin{equation}\label{eq:counterexample-commuting}
[L(f),g]+[f,L(g)]=0,~\text{whenever}~[f,g]=0.
\end{equation}

Finally, we prove the statement (i).
Set $F_1=F_2=L$ and $F_3=F_4=-L$.
If $fg=gf=0$, then $[f,g]=0$, and hence
\[
\begin{aligned}
&F_1(f)g+fF_2(g)+F_3(g)f+gF_4(f)\\
&=L(f)g+fL(g)-L(g)f-gL(f)\\
&=[L(f),g]+[f,L(g)]=0
\end{aligned}
\]
by \eqref{eq:counterexample-commuting}. By our assumption, there exist $C\in\I$,
a derivation $D:\I\to \I$ and a central-valued map $\Omega: \I\to Z(\I)$ such that
\begin{equation}\label{eq:L-assumed-standard}
L(f)=Cf+D(f)+\Omega(f)
\end{equation}
for all $f\in\I$. Writing $\one=e_z+\sum_i p_i$, we have from the construction of $L$ that
$L(\one)=-\sum_i t_ip_i+\sum_i t_ip_i=0$. Taking $f=\one$ in \eqref{eq:L-assumed-standard},
we obtain $C=-\Omega(\one)\in Z(\I)$. Since the poset $X$ is connected, it follows from
Proposition~\ref{prop:center} that $C=c\one$ for some $c\in R$.
For any $1\leq i\leq m$, let $x\in X_i$. Taking $f=e_x$ in \eqref{eq:L-assumed-standard}, we have
\begin{equation}\label{eq:diagonal-contradiction}
t_i e_x=ce_x+D(e_x)+\omega_x\one,
\end{equation}
where $\omega_x\in R$ with $\omega_x\one=\Omega(e_x)$. Let $y\neq x$.
Multiplying $e_y$ in \eqref{eq:diagonal-contradiction} from both sides, we can get from Lemma~\ref{lem:derivation-diagonal}
that $\omega_x=0$. Furthermore, there is $t_i=c$, since $D(e_x)(x,x)=0$ also by Lemma~\ref{lem:derivation-diagonal}.
In other words, all the $t_i$'s are equal to each other, contradicting the
choice of the $t_i$ for $1\leq i\leq m$. Thus (i) holds.
\end{proof}

\end{document}
